\chapter{Market and Technology Survey}
\label{chap:2}

\section{Survey Method}
\label{sec:2.1}

Five commercial products were surveyed from August to September 2026. A product qualified on three conditions.

First, it claims to observe a user performing an exercise and give feedback on how the movement was performed. Products that only count repetitions, only prescribe a programme, or only track heart rate were excluded.

Second, it addresses general strength or bodyweight training. Products confined to a single specialised discipline were excluded, because their movement vocabulary and their failure modes do not compare.

Third, enough is published about it to place it on every criterion below. A product described only in summary would enter the comparison with blank cells, and a table with blanks invites conclusions it cannot support.

\subsection{Comparison criteria}
\label{sec:2.1.1}

{\small
\setlength{\LTpre}{6pt}\setlength{\LTpost}{8pt}
\begin{longtable}{>{\raggedright\arraybackslash}p{\dimexpr 0.417\textwidth-2\tabcolsep\relax}>{\raggedright\arraybackslash}p{\dimexpr 0.583\textwidth-2\tabcolsep\relax}}
\caption{The criteria every product and system in this chapter is compared against}\\
\toprule
Criterion & Why it is compared \\
\midrule
\endfirsthead
\toprule
Criterion & Why it is compared \\
\midrule
\endhead
\bottomrule
\endfoot
\bottomrule
\endlastfoot
Hardware required & Decides the addressable market and the setup friction \\
Inference location & Drives latency, privacy exposure and running cost \\
Feedback timing & During the repetition, after the set, or later \\
Feedback modality & Whether the user must look at a screen \\
Explains the cause & Whether the product says why, not only what \\
Published evaluation & Whether an accuracy claim can be checked \\
\end{longtable}
}

The survey establishes what each product claims and is reported to do. It does not establish how well any of them works, and 2.4 explains why no survey could.

\section{Commercial Products}
\label{sec:2.2}

{\footnotesize
\setlength{\LTpre}{6pt}\setlength{\LTpost}{8pt}
\begin{longtable}{>{\raggedright\arraybackslash}p{\dimexpr 0.177\textwidth-2\tabcolsep\relax}>{\raggedright\arraybackslash}p{\dimexpr 0.197\textwidth-2\tabcolsep\relax}>{\raggedright\arraybackslash}p{\dimexpr 0.198\textwidth-2\tabcolsep\relax}>{\raggedright\arraybackslash}p{\dimexpr 0.265\textwidth-2\tabcolsep\relax}>{\raggedright\arraybackslash}p{\dimexpr 0.163\textwidth-2\tabcolsep\relax}}
\caption{Five commercial products, and what each asks of the user}\\
\toprule
Product & Hardware and inference & Feedback & Coverage & Price \\
\midrule
\endfirsthead
\toprule
Product & Hardware and inference & Feedback & Coverage & Price \\
\midrule
\endhead
\bottomrule
\endfoot
\bottomrule
\endlastfoot
Kemtai & Webcam in a browser, on the device & Real-time, audio and visual & 2,000+ rehab and fitness movements & Not published \\
Zing Coach & Phone camera, cloud & Real-time, visual prompts & Bodyweight and machines & About \$18.99 per month \\
Ai\allowbreak{}KYNETIX & Phone camera, cloud & Real-time, visual overlay & Barbell lifts & \$12.99 per month \\
Form\allowbreak{}Check AI & Phone camera, cloud & About 15 s after the set, visual & 30+ movements & About \$12 to \$13 per month \\
SensAI & Phone camera, cloud & Asynchronous, conversational text & Programme-dependent & \$6.99 per month \\
\end{longtable}
}

Prices are as listed at the survey date and are not independently verified. The cloud product rows draw on a published category review (SensAI product review, 2026)\cite{sensai2026review}, which is itself published by one of the products listed.

Three patterns are worth drawing out.

Four of the five run inference in the cloud. That places a network round trip inside the feedback path, and it gives the vendor a cost for every repetition observed, which has to be recovered from the subscription.

Only one of the five speaks. The other four require the user to watch a screen while training, which is not possible at the bottom of a squat. The one that does speak needs a webcam and a browser, not a phone.

Two of the five have left the category without saying so. FormCheck AI returns its analysis about fifteen seconds after the set and SensAI is asynchronous altogether, so neither can change the repetition it is describing. At that point form feedback has stopped being coaching and become review.

\section{Research Background}
\label{sec:2.3}

\subsection{Pose estimation on a mobile device}
\label{sec:2.3.1}

{\footnotesize
\setlength{\LTpre}{6pt}\setlength{\LTpost}{8pt}
\begin{longtable}{>{\raggedright\arraybackslash}p{\dimexpr 0.193\textwidth-2\tabcolsep\relax}>{\raggedright\arraybackslash}p{\dimexpr 0.272\textwidth-2\tabcolsep\relax}>{\raggedright\arraybackslash}p{\dimexpr 0.288\textwidth-2\tabcolsep\relax}>{\raggedright\arraybackslash}p{\dimexpr 0.247\textwidth-2\tabcolsep\relax}}
\caption{Mobile pose estimation engines, with the figures their authors publish}\\
\toprule
Engine & Keypoints and topology & Published figure & Source \\
\midrule
\endfirsthead
\toprule
Engine & Keypoints and topology & Published figure & Source \\
\midrule
\endhead
\bottomrule
\endfoot
\bottomrule
\endlastfoot
Blaze\allowbreak{}Pose & 33, including heel and foot index & Over 30 fps on a Pixel 2, single person & (Bazarevsky et al., 2020)\cite{bazarevsky2020blazepose} \\
Media\allowbreak{}Pipe Pose Landmarker & 33, the shipped form of Blaze\allowbreak{}Pose & Lite, Full and Heavy variants & (Google Media\allowbreak{}Pipe documentation)\cite{google2024mediapipe} \\
MoveNet & 17, COCO topology & Pixel 5 latency and COCO mAP published & (Tensor\allowbreak{}Flow Blog, MoveNet, 2021)\cite{tensorflow2021movenet} \\
RTMPose & Configurable, COCO topology by default & Real-time multi-person & (Jiang et al., 2023)\cite{jiang2023rtmpose} \\
\end{longtable}
}

These figures establish that real-time pose estimation on a commodity phone is feasible. They do not establish performance on this project's device, and none of them is carried forward as though it did.

These four differ in output topology as well as in speed. The COCO seventeen-keypoint layout stops at the ankle, with no heel and no toe, while BlazePose carries thirty-three points including both. Topology is the stronger constraint where a fault depends on a landmark the layout does not contain, because no amount of speed recovers a point the model never emits. Every fault this project retains is expressible in either layout, so the choice between them rests on measurement rather than on topology.

\subsection{Movement assessment from skeleton sequences}
\label{sec:2.3.2}

{\small
\setlength{\LTpre}{6pt}\setlength{\LTpost}{8pt}
\begin{longtable}{>{\raggedright\arraybackslash}p{\dimexpr 0.259\textwidth-2\tabcolsep\relax}>{\raggedright\arraybackslash}p{\dimexpr 0.413\textwidth-2\tabcolsep\relax}>{\raggedright\arraybackslash}p{\dimexpr 0.328\textwidth-2\tabcolsep\relax}}
\caption{Published systems that assess movement from skeleton sequences}\\
\toprule
System & What it does & Data it used \\
\midrule
\endfirsthead
\toprule
System & What it does & Data it used \\
\midrule
\endhead
\bottomrule
\endfoot
\bottomrule
\endlastfoot
Pose Trainer (Chen and Yang, Stanford, 2020)\cite{chen2020posetrainer} & Form recommendations for four exercises, evaluating geometric heuristics against learned classifiers & Over 100 videos the authors recorded, correct and deliberately incorrect \\
Fitness-AQA (Parmar, Gharat and Rhodin, ECCV 2022)\cite{parmar2022fitnessaqa} & Self-supervised representations for workout form assessment & Real gym video, form errors annotated by professional trainers \\
EC3D (Zhao et al., ACCV 2022)\cite{zhao2022ec3d} & Graph convolutional network over pose sequences, identifying the mistake and proposing a correction & Squat, lunge and plank, correct and with deliberate mistakes, four subjects \\
\end{longtable}
}

All three analyse recorded sequences after the fact. Offline, the whole repetition is available before a verdict is required. Live, the system has to commit while the movement is still happening and while a cue can still change it. That is a different problem, and the accuracy figures these systems report do not transfer to it.

\subsection{Where the fault definitions come from}
\label{sec:2.3.3}

This project detects six technique faults across its three exercises: in the squat, shallow depth and forward lean; in the push-up, the hips sagging and the hips piking; in the plank, the hips sagging and the hips piking. Four further faults were considered and excluded, each with its reason recorded in the specification file that defines the exercise. Two are not observable in the single camera view those exercises declare, one describes a deviation a spoken cue cannot correct, and one has no source that defines it at all.

The set is adopted from published work rather than invented. The Functional Movement Screen supplies scored observational criteria for the Deep Squat and the Trunk Stability Push-up (Cook, Burton and Hoogenboom 2006)\cite{cook2006fms}, covering both squat faults and push-up hip sag. The plank endurance test termination criterion, consistent across protocols (Eimiller et al., 2025; Rodríguez-Perea et al., 2025)\cite{eimiller2025plank,rodriguezperea2025plank}, names the two failure modes as the hips lowering and the pelvis rising, which are the two plank faults. Those are test-termination rules rather than a taxonomy annotated by trainers, which makes them weaker evidence than the Functional Movement Screen, and the specification file says so. Fitness-AQA corroborates shallow depth in the squat as a fault that trainers annotate, though its exercises are barbell movements, which this project does not cover.

Five of the six faults carry a cited criterion on this basis. The one that does not is push-up piking, which rests on symmetry with a sourced fault, and that is an argument rather than a citation. It is marked as unsourced in the specification file rather than left to look settled.

What no source supplies is more consequential. Every source establishes that a deviation is a fault and gives an observable condition for recognising it. None gives the quantity at which the fault begins. The specifications therefore record the taxonomy source and the threshold source as separate fields, and the threshold field is empty throughout. No result from this project can therefore be reported without stating how its threshold was derived.

\section{Summary: What This Market Does Not Provide}
\label{sec:2.4}

Three questions were put to all five products.

{\small
\setlength{\LTpre}{6pt}\setlength{\LTpost}{8pt}
\begin{longtable}{>{\raggedright\arraybackslash}p{\dimexpr 0.500\textwidth-2\tabcolsep\relax}>{\raggedright\arraybackslash}p{\dimexpr 0.500\textwidth-2\tabcolsep\relax}}
\caption{Three questions put to all five products, and which of them answer each}\\
\toprule
Question & Products that answer it \\
\midrule
\endfirsthead
\toprule
Question & Products that answer it \\
\midrule
\endhead
\bottomrule
\endfoot
\bottomrule
\endlastfoot
Does it explain the cause of a fault, rather than only naming it? & None during a set. One explains in prose, but only asynchronously \\
Does it publish how its accuracy was evaluated? & None \\
Does it state how a new exercise is added? & None. Coverage is announced as a number, never as a mechanism \\
\end{longtable}
}

The second row matters most. Accuracy claims do appear in this market, but a figure published without the procedure that produced it, without the data it was measured on, and without any statement of who labelled that data cannot be checked, compared, or reproduced. A buyer therefore has no basis for preferring one product's accuracy over another's.

This is not an accusation. Publishing an evaluation earns a vendor nothing commercially and no regulation requires it. But it leaves a gap that a project judged on method rather than on sales is in an unusual position to close.

\subsection{Language}
\label{sec:2.4.1}

None of the five surveyed products offers coaching in Vietnamese. Two smaller applications list Vietnamese among their supported languages. A store listing describes the interface and the store page, not the spoken coaching, and the spoken cues those applications demonstrate are in English.

The finding this chapter records is therefore the narrow one: no product was found that delivers spoken coaching cues in Vietnamese during a set. For a product whose feedback path is audio by design, and whose users train in Vietnamese, that is not a cosmetic difference.

\subsection{Where FORMA sits}
\label{sec:2.4.2}

Four properties, each of which separates this project from all five products above. No product in the survey offers them as a set.

\textbf{A phone the user already owns, and the video never leaves it.} One product of the five keeps inference on the device, and it needs a webcam and a desktop browser. The other four send the user's video to a server, which puts a network round trip inside the feedback path and a cost per repetition behind the subscription. FORMA requires no account, no upload and no network.

\textbf{It speaks while the user is still moving.} A person at the bottom of a squat cannot read a screen, so four of the five are delivering their correction to someone who is not looking at it. Two of those four arrive after the set has finished, which is too late to change the repetition being described. Audio is not an enhancement of the interface here; during a set it is the interface.

\textbf{It is built to coach in Vietnamese.} No surveyed product does. Every spoken line is a pre-authored asset keyed by language in the specification files, so a Vietnamese cue set is a content addition and not a change to the feedback path. The bank currently holds English only, and nothing beyond that is claimed here.

\textbf{A new exercise is a specification file, not a release.} Three exercises run through one evaluator with no exercise-specific code anywhere in it, which is a passing test rather than a claim. No product in the survey states how an exercise is added at all; coverage is announced as a number and the mechanism behind it is never described.

Underneath these four is a quieter property: they can be checked. Every cue traces to a published fault taxonomy, and every figure this project reports names the procedure that produced it.
